% filas del capítulo 11: símbolo & significado & unidad \\ (4.ª columna, macro, opcional; toda fila termina en \\)
$a$, $\mathrm{df}$, $\mathrm{pv}$ & monto, factor de descuento y valor presente de un flujo del calendario, $\mathrm{pv} = a\,\mathrm{df}$ & moneda del flujo; adimensional; moneda \\
$a^{\MXN}$ & monto de un flujo del calendario en \MXN (FIX si es liquidado, \emph{spot} si es proyectado en \USD) & \MXN \\
$a_{\mathrm{conf}}$, $a_{\mathrm{der}}$, $\delta_{\mathrm{conf}}$ & monto confirmado, monto derivado y su diferencia $a_{\mathrm{conf}} - a_{\mathrm{der}}$ & moneda del flujo \\
$V^{\mathrm{p}}$ & PV propio de la operación (\texttt{pv\_mxn}) & \MXN \\
$V^{\mathrm{b}}_{\mathrm{csv}}$, $V^{\mathrm{b}}$ & PV del banco en su moneda (CSV) y convertido a \MXN & moneda; \MXN \\
$\Delta$, $\delta$ & diferencia $V^{\mathrm{p}} - V^{\mathrm{b}}$ y su razón al nocional en \MXN & \MXN; adimensional \\
$N^{\MXN}$ & nocional vigente de la operación en \MXN (en \USD, por el \emph{spot}) & \MXN \\
$\epsilon$ & tolerancia de la conciliación (fracción del nocional; por omisión $10^{-4}$) & adimensional \\
$d^{-}$, $d^{+}$ & día hábil (lunes a viernes) anterior y posterior a la fecha de valuación & fecha \\
$\mathcal{V}$ & conjunto de PV propios (en $d$ y en los días vecinos vivos) que fija el rango de fecha & \MXN \\
$g(x)$, $x^{*}$, $\hat{g}_n$ & PV propio con el mercado desplazado $x$ menos el del banco, su raíz y su pendiente por diferencia finita & \MXN; \pb o \%; \MXN por unidad \\
$h$, $\bar{x}$ & paso de la pendiente y cota de búsqueda de Newton (curva: $1~\pb$ y $1\,000~\pb$; \emph{spot}: $0.01~\%$ y $50~\%$) & \pb o \% \\
